% =============================================================================
% Section 03 — Unity of Source and Aggregate
% Book I — Ontology of Viability
% Chapter 03 — Unity and Plurality
%
% Purpose:
%   To distinguish unity attributable to a common source
%   from unity attributable to an aggregate, and to examine
%   their different relations to identity.
%
% Conceptual Policy:
%   1. Unity of source is not equivalent to unity of identity.
%   2. Unity of aggregate is not equivalent to unity of identity.
%   3. A common source does not automatically establish
%      numerical or ontological identity among its manifestations.
%   4. An aggregate does not automatically constitute
%      a single identity.
%   5. Neither source nor aggregate is assumed to be
%      the exclusive foundation of identity.
%   6. The relation between source, aggregate, and identity
%      remains open to ontological investigation.
%
% =============================================================================

\section{Unity of Source and Aggregate}
\label{sec:unity-of-source-and-aggregate}

% -----------------------------------------------------------------------------
\subsection{The Problem of Two Different Forms of Unity}

Unity may be attributed to an existent or to a plurality
of existents on different ontological grounds.

Two such grounds require particular distinction:
the unity of source and the unity of aggregate.

The unity of source concerns the relation of several
manifestations, entities, or determinations to a common
origin or ontological source.

The unity of aggregate concerns the relation through
which multiple components or elements constitute
an organized or bounded whole.

Although both may be described as forms of unity,
they do not express the same ontological relation.

A common source may be associated with multiple
manifestations without those manifestations necessarily
constituting one aggregate.

Conversely, an aggregate may consist of components
whose origins are different or whose relations to
a common source are not established.

Neither form of unity can therefore be substituted
for the other.

More importantly, neither establishes the unity
of identity without further ontological argument.

The present section examines these distinctions
and clarifies the limits of the inferences that
may be drawn from source and aggregate.

% -----------------------------------------------------------------------------
\subsection{The Meaning of Source}

The term \emph{source} designates an ontological point
of origin or grounding in relation to something
that originates, proceeds, or is manifested.

The concept is used here in an ontological rather
than merely chronological sense.

A source is not necessarily an earlier event in time,
a spatially separate object, or a hidden material
entity located behind what appears.

The relation between a source and its manifestations
may involve different ontological forms and cannot
be reduced in advance to temporal succession,
physical production, or spatial separation.

Likewise, the source must not be understood as
a mere conceptual label assigned to several
manifestations because they share certain features.

The ontological relation of source requires
independent examination.

To identify a common source is to establish
a particular relation among what is considered
to originate from or be grounded in that source.

It does not, by itself, establish that all such
manifestations are numerically identical,
that they constitute one aggregate,
or that they possess one identity.

The source relation must therefore be distinguished
from the identity relation.

% -----------------------------------------------------------------------------
\subsection{Unity of Source}

Unity of source refers to the condition in which
multiple manifestations or entities are related
to one and the same source.

This unity concerns the origin or grounding
of the related manifestations.

It does not necessarily concern their internal
composition, their structural organization,
or the number of identities involved.

Several manifestations may be related to one source
while remaining distinct manifestations.

The unity of their source does not erase their
differences, nor does it automatically establish
that those differences are merely variations
within one identity.

Similarly, the existence of multiple manifestations
does not automatically imply that the source itself
is plural.

The relation between one source and multiple
manifestations therefore permits questions that
cannot be answered through numerical counting alone.

For example, it remains necessary to examine:

\begin{enumerate}
    \item Whether the manifestations are directly
    related to the source or through intermediate
    ontological relations.

    \item Whether the source relation is constitutive,
    generative, grounding, or of another kind.

    \item Whether the manifestations preserve
    a relevant ontological relation to their source
    throughout their transformations.

    \item Whether the common source has any
    determinate implications for the identity
    of the manifestations.
\end{enumerate}

These questions clarify that unity of source
is a specific relation and not a complete account
of the unity of what originates from it.

% -----------------------------------------------------------------------------
\subsection{The Meaning of Aggregate}

An \emph{aggregate} is understood here as a plurality
of elements or components considered together
under some form of relation, organization,
association, or boundedness.

The concept does not presuppose that the aggregate
possesses one identity.

Nor does it presuppose that every aggregate is
merely an arbitrary collection without ontological
unity.

An aggregate may exhibit different degrees and
forms of organization.

Its components may be spatially adjacent,
structurally connected, materially integrated,
functionally coordinated, or related through
other conditions.

These relations may provide grounds for attributing
unity to the aggregate in a specified respect.

However, the presence of such unity does not
automatically determine the ontological status
of the aggregate as an identity.

The aggregate may be considered unified as
a composition, system, configuration, or organized
plurality while the question of its identity remains
distinct.

Consequently, the unity of an aggregate must be
examined in relation to the specific basis on which
its components are considered together.

% -----------------------------------------------------------------------------
\subsection{Unity of Aggregate}

Unity of aggregate refers to the unity attributable
to a plurality of elements considered as a whole.

This unity may arise through relations among
components, their organization, their integration,
or the conditions that delimit the aggregate.

The relevant unity is therefore not necessarily
a unity of origin.

An aggregate may include components that originate
from different sources.

Their composition may nevertheless constitute
a unified whole in a particular ontological respect.

Conversely, components that share a common source
may remain separate and may never constitute
one aggregate.

The distinction can be expressed schematically:

\[
\text{Unity of Source}
\neq
\text{Unity of Aggregate}
\]

This distinction does not imply that source and
aggregate are unrelated.

A source may be relevant to the formation of
an aggregate, and an aggregate may preserve
or express relations to a common source.

However, such relations require examination
rather than assumption.

The unity of an aggregate is consequently determined
by the relevant relations among its elements,
whereas unity of source concerns the relation
of manifestations or entities to their origin
or grounding.

Neither definition entails the other.

% -----------------------------------------------------------------------------
\subsection{Source and Aggregate as Distinct Ontological Relations}

The distinction between source and aggregate
can be clarified by considering the different
directions of the relations involved.

In the case of source, the inquiry proceeds
from manifestations or entities toward the
ontological ground or origin with which they
are related.

In the case of aggregate, the inquiry concerns
the relations among elements through which
they constitute a whole.

These are different relational orientations.

A common source may stand in relation to several
manifestations without those manifestations
being components of one aggregate.

Likewise, several components may constitute
an aggregate without sharing a common source.

The distinction can be represented as follows:

\[
\begin{aligned}
\text{Source Relation:}\quad
& S \longrightarrow \{M_1,M_2,\ldots,M_n\}
\\[6pt]
\text{Aggregate Relation:}\quad
& \{C_1,C_2,\ldots,C_n\}
\longrightarrow A
\end{aligned}
\]

Here, \(S\) designates a source,
\(M_i\) designates a manifestation,
\(C_i\) designates a component,
and \(A\) designates an aggregate.

These expressions are schematic.
They do not establish that every source produces
multiple manifestations or that every aggregate
has a single determinate source.

Their purpose is to distinguish two different
ontological relations.

The source relation concerns the relation
between a source and what is related to it.
The aggregate relation concerns the relation
among components considered as a whole.

The distinction is essential because the unity
of the relational subject in one case cannot
be transferred automatically to the subject
of the other.

% -----------------------------------------------------------------------------
\subsection{Why a Common Source Does Not Establish One Identity}

A common source may be relevant to the identity
of its manifestations.

However, the existence of a common source does
not by itself establish that all manifestations
belong to one identity.

Such an inference would require the additional
assumption that sharing a source is sufficient
for identity unity.

That assumption cannot be accepted without
independent justification.

Several manifestations may be related to one source
while differing in their ontological constitution,
relations, determinations, or modes of existence.

Whether these differences are compatible with
one identity or involve a plurality of identities
is a separate question.

The source relation alone does not settle it.

For example, the fact that several manifestations
originate from one source does not establish
that they are merely numerically distinct appearances
of one identity.

Nor does it establish the opposite conclusion:
that each manifestation necessarily constitutes
a separate identity.

Both interpretations exceed what the source relation
alone can demonstrate.

The correct ontological procedure is therefore
to recognize the common source where it is established
while independently examining the identity
of the manifestations.

This preserves the significance of source without
reducing identity to origin.

% -----------------------------------------------------------------------------
\subsection{Why an Aggregate Does Not Establish One Identity}

The unity of an aggregate also fails to establish,
by itself, the unity of identity.

An aggregate may possess internal organization,
functional coordination, material integration,
or structural coherence.

These conditions may justify describing it as
one aggregate in a specified respect.

Nevertheless, the unity of its composition
does not automatically determine whether it
possesses one identity, multiple identities,
or another ontological relation between
its constituents and its whole.

The inference from aggregate unity to identity unity
would require an additional account of the relation
between composition and identity.

Without such an account, the inference remains
unjustified.

The same restriction applies in the opposite direction.

The presence of multiple components does not
necessarily imply that the aggregate lacks
a unified identity.

Plurality of components and unity of identity
are not mutually exclusive by definition.

The question must be examined without reducing
identity either to the aggregate or to its components.

This distinction provides the conceptual basis
for the subsequent examination of plurality
and composition.

% -----------------------------------------------------------------------------
\subsection{The Independence of Source, Aggregate, and Identity}

Source, aggregate, and identity must be treated
as distinct ontological concepts.

Their relations may be significant, but their
conceptual distinction must be preserved.

The following table summarizes their respective
questions and the limits of the inferences
that may be drawn from them.


The distinctions summarized above do not imply
that source, aggregate, and identity exist in
complete ontological isolation from one another.

A source may participate in the constitution
of an aggregate.
An aggregate may be relevant to the manifestation
of an identity.
An identity may involve relations to a source
and to a plurality of components.

What must be rejected is not their possible
interdependence, but the reduction of one
to another.

A theory of identity must therefore explain
how source and aggregate are ontologically
relevant without assuming that either one
is sufficient to constitute or determine identity.

% -----------------------------------------------------------------------------
\subsection{Possible Relations Between Source and Aggregate}

Although source and aggregate are distinct,
their relations may assume different forms.

A common source may give rise to several entities
that later become related within an aggregate.

An aggregate may contain components originating
from different sources.

A source may be associated with manifestations
that remain separate and do not constitute
a single aggregate.

An aggregate may also undergo changes in its
components while retaining some of the relations
through which it is considered unified.

These possibilities demonstrate that source
and aggregate do not stand in a single necessary
relation.

At least four configurations can be distinguished:

\begin{enumerate}
    \item One source related to several manifestations
    that do not constitute one aggregate.

    \item Several sources associated with components
    that constitute one aggregate.

    \item One source related to manifestations
    that also participate in one aggregate.

    \item Multiple sources and multiple aggregates
    connected through further ontological relations.
\end{enumerate}

These configurations are analytical possibilities.
They are not intended as an exhaustive classification
of every possible relation between sources
and aggregates.

Their significance is that no single relation
between source and aggregate can be assumed
as universally necessary.

The unity of source and the unity of aggregate
must therefore be investigated according to
the conditions under which each is established.

% -----------------------------------------------------------------------------
\subsection{The Limits of Origin-Based and Composition-Based Explanations}

The distinction developed in this section
reveals two possible reductions that must be avoided.

The first is an origin-based reduction of identity.

According to this reduction, the identity
of an existent would be determined entirely
by its source or origin.

Such an account risks treating the relation
of origin as sufficient to settle all questions
of identity.

However, the fact that two manifestations share
a source does not by itself establish that
they possess one identity.

The second is a composition-based reduction
of identity.

According to this reduction, the identity
of an existent would be determined entirely
by the organization or unity of its components.

Such an account risks treating the aggregate
as the exclusive foundation of identity.

However, the fact that components form a unified
aggregate does not by itself establish that
the aggregate possesses one identity.

Both reductions confuse a relevant ontological
relation with a complete account of identity.

Rejecting these reductions does not mean that
origin and composition are irrelevant.
It means that their relevance must be established
without allowing either to replace the ontological
question of identity.

% -----------------------------------------------------------------------------
\subsection{The Relation to Unity and Plurality}

The distinction between source and aggregate
also clarifies the broader problem of unity
and plurality.

A plurality of manifestations may be unified
with respect to source while remaining plural
with respect to manifestation.

A plurality of components may be unified
with respect to aggregate while remaining plural
with respect to composition.

Neither case determines, without further argument,
whether unity or plurality should be attributed
to identity.

The subject and respect of each claim must
therefore remain explicit.

It is possible to speak of:

\begin{itemize}
    \item Unity of source with plurality of manifestations.
    \item Unity of aggregate with plurality of components.
    \item Unity of identity with plurality of determinations.
    \item Plurality of identities within a broader
    aggregate or relational configuration.
\end{itemize}

These descriptions concern different ontological
respects and should not be treated as contradictory
merely because they involve unity and plurality
simultaneously.

The central requirement is to determine what
is unified, what is plural, and in which respect
the relevant determination applies.

This requirement prevents the assumption that
the recognition of one form of unity excludes
every possible form of plurality.

% -----------------------------------------------------------------------------
\subsection{Conclusion: Source and Aggregate Do Not Settle Identity}

Unity of source and unity of aggregate are
distinct forms of ontological unity.

The former concerns the relation of manifestations
or entities to a common source.
The latter concerns the relation among components
considered as an organized or bounded whole.

Neither form is equivalent to unity of identity.

A common source does not automatically establish
that its manifestations belong to one identity.
Likewise, a unified aggregate does not automatically
establish that its components constitute one identity.

At the same time, neither source nor aggregate
is excluded from contributing to the ontological
conditions under which identity is constituted
or manifested.

Their relevance must be examined without reduction.

The distinction established here provides
the basis for the next stage of the inquiry:
the examination of plurality of identities.

That examination must determine whether and under
what conditions an existent or a broader ontological
configuration may involve more than one identity,
without deriving such plurality merely from
the multiplicity of manifestations or components.

% =============================================================================
% End of Section 03
% =============================================================================